% !TEX program = xelatex
% Compile from within slides/: xelatex -interaction=nonstopmode results_1a_1b_2a_presentation.tex
%
% Comprehensive presentation covering Study 1a (9 LLMs, zh/en hate detection),
% Study 1b (9 LLMs, 3 scales paired gender items), and Study 2a (779 human
% rating sessions from 712 unique participants, 3 scales paired gender items
% + ICC similarity + cross-scale).
% Structure per frame: hypothesis -> method -> result -> verdict.
% Style reference: 1b_groupswap_bias_presentation.tex.

\documentclass[aspectratio=169,professionalfonts]{beamer}
\usetheme{Madrid}
\usecolortheme{default}
\setbeamertemplate{navigation symbols}{}
\setbeamertemplate{footline}[frame number]

\usepackage{fontspec}
\usepackage{xeCJK}
\setmainfont{Latin Modern Roman}
\setsansfont{Latin Modern Sans}
\setmonofont{Latin Modern Mono}
\setCJKmainfont{Noto Serif CJK SC}
\setCJKsansfont{Noto Sans CJK SC}
\setCJKmonofont{Noto Sans Mono CJK SC}
\IfFontExistsTF{Noto Serif CJK SC}{}{
  \setCJKmainfont{Source Han Serif SC}
}
\IfFontExistsTF{Noto Sans CJK SC}{}{
  \setCJKsansfont{Source Han Sans SC}
}
\IfFontExistsTF{Noto Sans Mono CJK SC}{}{
  \setCJKmonofont{Source Han Sans SC}
}

\usepackage{amsmath}
\usepackage{booktabs}
\usepackage{array}
\usepackage{graphicx}
\usepackage{tabularx}
\usepackage{multirow}
\usepackage{ragged2e}
\usepackage{xcolor}

% Paths relative to the slides/ directory (same convention as the paper tex).
\newcommand{\artifacts}{../artifacts}
\newcommand{\oneb}{../outputs/group_swap_1b/1b_groupswap_demensionsentence/analysis_1b/figures}
\newcommand{\humanfig}{../human/outputs/final_clean_long_gender_difference_analysis/figures}
\newcommand{\humanpgfig}{../human/outputs/final_clean_long_participant_gender_combined_analysis/figures}

% Shorthand blocks
\newcommand{\Hyp}[2]{\textbf{#1}~#2}

\title[研究 1a / 1b / 2a 综合结果汇报]{大语言模型与人类在中文仇恨言论识别与性别偏差评分中的一致性与异质性\\[0.3em]\large 研究 1a、1b、2a 综合结果汇报}
\author{鲍璇}
\date{\today}

\begin{document}

% ===========================================================================
\begin{frame}
  \titlepage
\end{frame}

\begin{frame}{报告结构}
  \begin{enumerate}
    \item 总览：三项研究的问题链与共用语料
    \item 研究 1a：9 个 LLM（Large Language Model，大语言模型）中英文仇恨识别（H1a.1 vs H1a.2）
    \item 研究 1b：9 个 LLM 性别配对评分（H1b.1、H1b.2、H1b.3）
    \item 研究 2a：779 个人类评分 session 性别配对评分（H2a.1--H2a.4；与 1b 配对）
    \item 研究 2a$\cdot$ICC（Intraclass Correlation Coefficient，类内相关系数）：人机偏差模式题项级相似度（H2a.ICC.1--.3）
    \item 研究 2a$\cdot$CS（Cross-Scale similarity，跨尺度相似度）：人机跨尺度稳健性比较（H2a.CS.1--.2）
    \item 综合结论与研究局限
  \end{enumerate}
\end{frame}

\begin{frame}{统计量速览（正文指标速览）}
  \footnotesize
  \setlength{\tabcolsep}{4pt}
  \begin{tabular}{p{0.22\textwidth}p{0.73\textwidth}}
    \toprule
    符号 / 检验 & 读法与判读准则 \\
    \midrule
    $\Delta_{F-M}$ & 同一对语句中\,女人版\,减\,男人版\,的评分差；$>0$ 即女人版被评得更严重。 \\
    $d_z$         & 配对样本效应量；Cohen 约定：$0.2$ 小、$0.5$ 中、$0.8$ 大。本研究 $d_z$ 区间 $0.29$--$1.02$。 \\
    Spearman $\rho_s$ & 两个排序之间的相关（$-1$ 到 $+1$）；越接近 $+1$ 越一致。 \\
    方差比 $\frac{\text{Var}_{zh}}{\text{Var}_{en}}$ & 中文 vs 英文模型分数的\emph{离散度}之比；$\gg 1$ 即中文端被拉得更开。 \\
    Bartlett / Levene & 方差同质性检验，$p<.05$ 即拒绝"两组方差相等"。 \\
    FDR (False Discovery Rate，Benjamini--Hochberg) $q$ & 多重检验校正后的显著性阈值；$q<.05$ 即在控制错误发现率下仍显著。 \\
    Bootstrap 95\% CI (Confidence Interval) & 有放回重抽样 $B=2000$ 次得到的置信区间；不含 $0$ 即与零效应显著不同。 \\
    ICC(2,1) & Intraclass Correlation Coefficient（类内相关系数）；双向随机、单次测量的绝对一致性；Koo \& Li (2016) 分档：$<.50$ poor、$.50$--$.75$ moderate、$.75$--$.90$ good、$>.90$ excellent。 \\
    \bottomrule
  \end{tabular}

  \vspace{0.3em}
  \small 下文所有数字都可以回到这张表对齐解读。
\end{frame}

\begin{frame}{三项研究的问题链}
  \begin{columns}[T,onlytextwidth]
    \column{0.32\textwidth}
      \begin{block}{研究 1a}
        \small 9 个 LLM 能否在广泛的中英文仇恨识别任务上表现出\emph{跨语言稳定}的能力？若不然，差距来自哪里？在特定的性别仇恨上又是如何表现？
      \end{block}
    \column{0.32\textwidth}
      \begin{block}{研究 1b}
        \small 仅替换性别指称词，9个LLM、3种计分方式对\emph{女人版}与\emph{男人版}攻击性评分是否存在系统性不同？
      \end{block}
    \column{0.32\textwidth}
      \begin{block}{研究 2a}
        \small 中文母语人类对1b同样的内容（仅替换性别指代词）评分，人类是否呈现同LLM相似的方向偏差？男女是否有模式差异？哪些LLM模式与人类最接近？
      \end{block}
  \end{columns}

  \vspace{0.5em}
  \begin{itemize}
    \item \textbf{研究1b、2a共用语料}：371 对中文同质配对语句（minimal pairs）；3 层编码框架（10 一级 / 71 二级 / 183 三级攻击维度）。
    \item \textbf{共用评分尺度}：3 点 / 7 点 Likert / 0--100 滑块（三条平行测量）。
    \item \textbf{9 个 LLM}：Claude Opus 4.5、GPT-5.1、Gemma 4 31B、Llama 4 Maverick、Qwen 2.5 72B、GLM 4.6、Kimi K2、DeepSeek R1、DeepSeek V3.2。
  \end{itemize}
\end{frame}

% ===========================================================================
\section{研究 1a：9 个 LLM 中英文仇恨识别}

\begin{frame}{研究 1a：研究问题}
  \begin{block}{核心问题}
    \small 9 个 LLM 在中、英文仇恨言论识别上表现出可观察的差距。该差距是由\textbf{任务定义的跨语言分歧}造成，还是由\textbf{训练语料对中文的暴露不均}造成？
  \end{block}

  \vspace{0.4em}
  \begin{itemize}
    \item 9 个 LLM $\times$ 2 语言（中 / 英）$\times$ 2 提示（zero-shot / CoT，Chain-of-Thought，链式思考）$=$ 36 个评估单元。
    \item 指标：F1（分类准确度）+ Brier（概率校准）。
    \item 两条可证伪、互斥的假设作为框架。
  \end{itemize}
\end{frame}

\begin{frame}{研究 1a：两条竞争假设}
  \begin{block}{H1a.1 任务定义分歧假设}
    \small 若差距源于中英标注规范与文化边界的分歧，差距应\emph{对所有模型为近似常数}——两语言 F1／Brier 分布仅发生整体平移，\emph{跨模型方差大致相当}、排序稳定。
  \end{block}

  \begin{block}{H1a.2 语言覆盖不均假设}
    \small 若差距源于训练语料对中文的覆盖不均（各厂商中文语料差异悬殊），则\emph{中文端模型间的方差应显著大于英文端}。H1a.2 的关键特征不是"排名谁变了"，而是"中文上模型之间的差距更大"。
  \end{block}

  \vspace{0.4em}
  两种假设的关键分界：\textbf{等方差（H1a.1）vs.\ 方差膨胀（H1a.2）}。
\end{frame}

\begin{frame}{研究 1a：散点与斜率图——中文端模型间差异更大}
  \begin{columns}[T,onlytextwidth]
    \column{0.5\textwidth}
      \centering
      \includegraphics[width=\linewidth,height=0.70\textheight,keepaspectratio]{\artifacts/validated_csv_hate_f1_brier_2x2_scatter.pdf}
    \column{0.5\textwidth}
      \centering
      \includegraphics[width=\linewidth,height=0.70\textheight,keepaspectratio]{\artifacts/validated_csv_hate_f1_brier_language_slope.pdf}
  \end{columns}
  \vspace{0.3em}
  \small\centering 左：4 个面板的 F1--Brier 联合散点；右：每个模型从中文到英文的斜率图（菱形 = 面板均值 $\pm$ SD）。
\end{frame}

\begin{frame}{研究 1a：中文面板 SD 是英文的 3--4 倍——方差已显现不对称}
  \footnotesize
  \setlength{\tabcolsep}{4pt}
  \begin{table}
    \centering
    \caption{\small 四个面板的描述性统计与最佳模型}
    \begin{tabular}{llcccc}
      \toprule
      语言 & 提示 & F1 均值 $\pm$ SD & 最佳 F1 & Brier 均值 $\pm$ SD & 最佳 Brier \\
      \midrule
      中 & zero-shot & $0.608 \pm 0.085$ & $0.751$ (Qwen)  & $0.317 \pm 0.041$ & $0.240$ (Qwen) \\
      中 & CoT       & $0.602 \pm 0.077$ & $0.736$ (Llama) & $0.323 \pm 0.039$ & $0.243$ (Llama) \\
      英 & zero-shot & $0.654 \pm 0.024$ & $0.698$ (Llama) & $0.282 \pm 0.013$ & $0.254$ (Qwen) \\
      英 & CoT       & $0.652 \pm 0.020$ & $0.699$ (Llama) & $0.280 \pm 0.010$ & $0.256$ (Llama) \\
      \bottomrule
    \end{tabular}
  \end{table}

  \vspace{0.2em}
  \textbf{一眼可见}：中文面板的 SD（0.077--0.085）约为英文面板 SD（0.020--0.024）的 \textbf{3--4 倍}——方差显著不对称。
\end{frame}

\begin{frame}{研究 1a：F1 排序跨语言基本保留，但方差比 10--15 倍——拒绝等方差}
  \footnotesize
  \setlength{\tabcolsep}{4pt}
  \begin{table}
    \centering
    \caption{\small 跨语言 Spearman $\rho_s$ 与方差同质性检验}
    \begin{tabular}{llcccc}
      \toprule
      指标 & 提示 & $\rho_s$（$p$） & 方差比 $\frac{\text{Var}_{zh}}{\text{Var}_{en}}$ & Bartlett（$p$） & Levene（$p$） \\
      \midrule
      F1    & zero-shot & $+0.883$（.002） & $12.68$ & $9.83$（.002） & $6.17$（.024） \\
      F1    & CoT       & $+0.750$（.020） & $15.00$ & $10.93$（$<$.001） & $9.40$（.007） \\
      Brier & zero-shot & $+0.383$（.309） & $10.70$ & $8.75$（.003） & $5.57$（.031） \\
      Brier & CoT       & $+0.283$（.460） & $15.10$ & $10.97$（$<$.001） & $5.32$（.035） \\
      \bottomrule
    \end{tabular}
  \end{table}

  \begin{itemize}
    \small
    \item F1 跨语言排序\emph{基本保留}（$\rho_s = 0.75$--$0.88$）。
    \item 方差比 $10$--$15$ 倍，Bartlett / Levene \emph{一致拒绝}等方差。
    \item Brier 排序在跨语言上已经失稳。
  \end{itemize}
\end{frame}

\begin{frame}{研究 1a：两个关键发现}
  \begin{alertblock}{发现一：均值差距不大，两种假设都能解释}
    \small $\Delta_\text{F1}(zh-en)$：zero-shot $M = -0.046$（SD $0.069$，7/9 为负）；CoT $M = -0.050$（SD $0.063$，7/9 为负）。$\Delta_\text{Brier}(zh-en)$：约 $+0.04$（8/9 为正）。均值差距不大，H1a.1 与 H1a.2 都能大致解释。
  \end{alertblock}

  \begin{alertblock}{发现二：中文端模型之间差异大得多——这才是关键}
    \small 中文 F1 的标准差约为英文的 $12.7$ 倍（zero-shot）至 $15.0$ 倍（CoT）。这不是所有模型在中文上统一变差，而是\emph{有的模型中文很强，有的中文很弱}——这正是 H1a.2 所预测的。
  \end{alertblock}
\end{frame}

\begin{frame}{研究 1a：Gemma 中文塌陷、Qwen 中文反向领先——方差膨胀有具体模型来源}
  \centering
  \includegraphics[width=0.80\textwidth,height=0.56\textheight,keepaspectratio]{\artifacts/validated_csv_hate_f1_2x2_bars.pdf}

  \footnotesize
  \begin{itemize}\setlength{\itemsep}{0.1em}
    \item \textbf{Gemma-4-31B} 中文塌陷：$\Delta_{zh-en} = -0.171$（zero-shot），下端离群。
    \item \textbf{Qwen-2.5-72B} 中文反向领先：$\Delta_{zh-en} = +0.066$，上端离群。
    \item CoT 并没有缩小各模型间的差异（方差比：zero-shot $12.7\times$ $\to$ CoT $15.0\times$）。
  \end{itemize}
\end{frame}

\begin{frame}{研究 1a：CoT 并非普适增益，其方向取决于模型对该语言的掌握程度}
  \centering
  \includegraphics[width=0.88\textwidth,height=0.60\textheight,keepaspectratio]{\artifacts/validated_csv_hate_f1_2x2_dumbbell.pdf}

  \vspace{0.1em}
  \footnotesize
  \begin{itemize}
    \item 每根水平线 = 同一模型从 zero-shot 到 CoT 的 F1 轨迹；按语言分两面板。
    \item \textbf{CoT 并非普适增益}：Qwen 在中文 $0.751 \to 0.639$；Llama 在中文 $0.641 \to 0.736$；英文面板变化很小。
    \item CoT 能否带来提升，取决于模型对这门语言的基础掌握程度——这与 H1a.1 预测的"整体平移"不符。
  \end{itemize}
\end{frame}

\begin{frame}{研究 1a：Brier 校准复现同一方差膨胀格局——中文端离散度 10--15 倍}
  \centering
  \includegraphics[width=0.84\textwidth,height=0.68\textheight,keepaspectratio]{\artifacts/validated_csv_hate_brier_2x2_bars.pdf}

  \vspace{0.1em}
  \footnotesize Brier 越低越好。中文端各模型之间的概率校准差异，约是英文端的 $10.7$--$15.1$ 倍——与 F1 的结果一致，中文端方差更大。
\end{frame}

\begin{frame}{研究 1a：Brier 的"校准增益"与 F1 的"分类增益"不同步——两者可解耦}
  \centering
  \includegraphics[width=0.88\textwidth,height=0.62\textheight,keepaspectratio]{\artifacts/validated_csv_hate_brier_2x2_dumbbell.pdf}

  \vspace{0.1em}
  \footnotesize
  \begin{itemize}
    \item \textbf{CoT 对 Brier 非单调}：Llama 中文 $0.300 \to 0.243$；Qwen 英文 $0.254 \to 0.289$。
    \item CoT 对概率校准的影响与对 F1 的影响并不同步——F1 和 Brier 的高低不总一致。
  \end{itemize}
\end{frame}

\begin{frame}{研究 1a：F1 榜首不一定是最佳——Llama 分类好却过度自信，Qwen 概率更贴近真实}
  \footnotesize
  \begin{table}
    \centering
    \setlength{\tabcolsep}{5pt}
    \begin{tabular}{lcc}
      \toprule
      模型 & F1 (英文 zero-shot) & Brier (英文 zero-shot) \\
      \midrule
      Llama-4-Maverick & \textbf{0.698} (榜首) & $0.297$（偏高） \\
      Qwen-2.5-72B     & $0.686$ (第二)        & \textbf{0.254} (最低) \\
      \bottomrule
    \end{tabular}
  \end{table}

  \vspace{0.4em}
  \begin{block}{关键观察}
    \small Llama 分类准确率最高，但概率输出偏向过度自信；Qwen 分类略逊但概率\emph{更贴近真实频率}。如果后续需要用概率来设定审核门槛（比如"低于多少置信度就交给人工判断"），F1 第一不一定是最合适的选择——\textbf{F1 和 Brier 可以相互独立}。
  \end{block}
\end{frame}

\begin{frame}{研究 1a：H1a.1 被证伪、H1a.2 被支持——差距的主要来源是中文语料暴露不均}
  \small
  \begin{table}
    \centering
    \setlength{\tabcolsep}{4pt}
    \begin{tabular}{lll}
      \toprule
      假设 & 关键预测 & 数据判决 \\
      \midrule
      H1a.1 任务定义分歧 & 等方差、排序稳定 & \textbf{证伪}（方差比 $12.7$--$15.0$；Bartlett／Levene 均拒绝） \\
      H1a.2 语言覆盖不均 & 中文方差 $\gg$ 英文方差 & \textbf{支持}（同上 + Gemma 塌陷、Qwen 反向领先） \\
      \bottomrule
    \end{tabular}
  \end{table}

  \vspace{0.4em}
  \begin{block}{一行结论}
    \small 中英性能差距的主要来源不是标注定义上的分歧，而是各模型中文训练语料的多寡不同；严格驳回 H1a.1 还需要跨标注体系的证据，但当前数据中方差的不对称已超出 H1a.1 所能预测的范围。
  \end{block}
\end{frame}

% --- 1a 辅助 / 过渡到 1b --------------------------------------------------
\begin{frame}{研究 1a 过渡 $\to$ 1b：HateXplain 上"针对男性"帖反被打更高——样本构成混淆 $>$ 真实偏差}
  \centering
  \includegraphics[width=0.80\textwidth,height=0.46\textheight,keepaspectratio]{\artifacts/placeholder_a_1a_gender_target_mean_score.pdf}

  \scriptsize
  \begin{itemize}\setlength{\itemsep}{0.1em}
    \item 按 HateXplain \texttt{target} 字段抽取"仅 Women"（$n=175$）与"仅 Men"（$n=36$）两子集；9 个 LLM 在 zero-shot / CoT 下\emph{均}把"针对男性"帖打得更高（total：$3.19$ vs $2.64$；$3.35$ vs $2.72$），与"女性方向偏差"预期相反。
    \item \textbf{并不等于 LLM 保护男性}：Men 样本仅 $36$ 条且集中在与族群交叉的极端仇恨语境；Women 样本覆盖更广——\emph{"性别"与"语境极端性"被混淆}。
    \item \textbf{样本构成混淆 $>$ 真实偏差}——只有剥离语义内容、只留性别指称的\emph{最小配对设计}才能分离评估偏差，这正是研究 1b 的动因。
  \end{itemize}
\end{frame}

% ===========================================================================
\section{研究 1b：9 个 LLM 性别配对评分}

\begin{frame}{研究 1b：研究问题与设计}
  \small 承接 1a 过渡页的结论——自然数据中"性别 $\times$ 语境极端性"的样本构成混淆掩盖了真实评估偏差；1b 以\textbf{最小配对设计}消除语义差异，只保留性别指称作为唯一操纵变量。

  \begin{block}{核心问题}
    \small 在同一语义模板下，仅替换性别指称词（女人版 $\leftrightarrow$ 男人版），9 个 LLM 对攻击性的判定是否系统性不同？
  \end{block}

  \begin{itemize}
    \item \textbf{最小配对设计}：371 对中文同质句；唯一操纵变量 = 性别指称词。
    \item \textbf{评分尺度}：3 点（0--2）、7 点 Likert（0--6）、0--100 滑块——三条平行测量。
    \item \textbf{分析单位}：每对的 $\Delta_{F-M} = $ 女人版评分 $-$ 男人版评分；在 371 对上汇总为配对样本 Cohen's $d_z$。
  \end{itemize}
\end{frame}

\begin{frame}{研究 1b：三条可证伪的假设}
  \footnotesize
  \begin{block}{H1b.1 评估偏差假设（含语义内容零假设）}
    \small 模型携带独立于语义的、指向女性方向的偏差；$\Delta_{F-M}$ 应系统性 $> 0$；若偏差广泛存在，则不同模型应在\emph{方向}而非\emph{效应量}上一致。\emph{零假设对立面}：若判定完全由句意决定，27 个模型-尺度单元的 $\Delta_{F-M}$ 分布应以 $0$ 为中心对称。
  \end{block}
  \begin{block}{H1b.2 偏差维度集中假设}
    \small 偏差集中在与女性身体化／性化刻板印象绑定的领域（性化攻击、外貌形象），10 个一级领域上 $d_z$ 呈集中而非铺开的分布。
  \end{block}
  \begin{block}{H1b.3 题项级跨尺度同构假设}
    \small 同一模型在三尺度下的题项级 $\Delta_i$ 共享同一潜结构；跨尺度 Spearman $\rho$ 应接近 $1$；应在聚合／维度／题项三层粒度上一致稳健。
  \end{block}
  \begin{block}{}
    \small \textbf{与 2a 的配对关系}：H1b.1 $\leftrightarrow$ H2a.1（F$>$M 主效应）、H1b.2 $\leftrightarrow$ H2a.2（维度集中）、H1b.3 $\leftrightarrow$ H2a.CS.1（题项级跨尺度同构）。
  \end{block}
\end{frame}

\begin{frame}{研究 1b：27/27 单元 $q<10^{-7}$，$d_z\in[0.29,1.02]$——H1b.1 强支持}
  \centering
  \includegraphics[width=\textwidth,height=0.80\textheight,keepaspectratio]{\artifacts/overall_dz_three_scale_forest.pdf}

  \vspace{0.1em}
  \footnotesize 每个面板对应一种评分尺度；每条横线为该模型 $d_z$ 的 95\% bootstrap CI；三个面板共享 x 轴。\textbf{全部 27 个模型-尺度组合 $q < 10^{-7}$}；$d_z$ 从 $0.29$ 到 $1.02$，模型间相差最多约 3.5 倍。
\end{frame}

\begin{frame}{研究 1b：持平带随尺度分辨率从 $80\%$ 收缩到 $38\%$——量级被更细地记录}
  \centering
  \includegraphics[width=\textwidth,height=0.66\textheight,keepaspectratio]{\artifacts/directionality_three_scale_grid.pdf}

  \scriptsize 红 = 女人版更高；灰 = 持平；蓝 = 男人版更高。\textbf{持平的比例随尺度粒度变细而下降}：9 模型中位 $80.1\% \to 62.0\% \to 37.7\%$（3 点 $\to$ 7 点 $\to$ 滑块）——\emph{不是偏差减少了}，而是更细的尺度把原本"持平"的小差异也记录下来了。
\end{frame}

\begin{frame}{研究 1b：27 单元完整统计（1/2）——所有模型三尺度均显著 $\Delta_{F-M}>0$}
  \scriptsize
  \setlength{\tabcolsep}{3pt}
  \renewcommand{\arraystretch}{0.9}
  \begin{table}
    \centering
    \begin{tabular}{llccccc}
      \toprule
      模型 & 尺度 & $\Delta_{F-M}$ & 95\% CI & $d_z$ & 女\,:\,平\,:\,男 & $q$ (FDR) \\
      \midrule
      \multirow{3}{*}{Claude Opus 4.5}
        & 3 点    & 0.199 & $[0.159, 0.240]$ & 0.498 & 74\,:\,297\,:\,0   & $1.8\!\times\!10^{-17}$ \\
        & 7 点    & 0.345 & $[0.294, 0.396]$ & 0.693 & 128\,:\,241\,:\,2  & $1.5\!\times\!10^{-27}$ \\
        & 滑块    & 3.553 & $[3.084, 4.046]$ & 0.748 & 183\,:\,184\,:\,4  & $1.4\!\times\!10^{-31}$ \\
      \midrule
      \multirow{3}{*}{DeepSeek R1}
        & 3 点    & 0.105 & $[0.067, 0.143]$ & 0.289 & 46\,:\,318\,:\,7   & $8.9\!\times\!10^{-8}$ \\
        & 7 点    & 0.243 & $[0.186, 0.302]$ & 0.422 & 114\,:\,233\,:\,24 & $8.3\!\times\!10^{-14}$ \\
        & 滑块    & 4.326 & $[3.464, 5.243]$ & 0.510 & 203\,:\,116\,:\,52 & $5.7\!\times\!10^{-23}$ \\
      \midrule
      \multirow{3}{*}{DeepSeek V3.2}
        & 3 点    & 0.124 & $[0.081, 0.167]$ & 0.289 & 60\,:\,297\,:\,14  & $8.9\!\times\!10^{-8}$ \\
        & 7 点    & 0.237 & $[0.173, 0.305]$ & 0.371 & 114\,:\,223\,:\,34 & $2.2\!\times\!10^{-11}$ \\
        & 滑块    & 3.889 & $[3.054, 4.736]$ & 0.460 & 208\,:\,96\,:\,67  & $7.8\!\times\!10^{-17}$ \\
      \midrule
      \multirow{3}{*}{Gemma 4 31B}
        & 3 点    & 0.143 & $[0.108, 0.181]$ & 0.408 & 53\,:\,318\,:\,0   & $6.0\!\times\!10^{-13}$ \\
        & 7 点    & 0.385 & $[0.334, 0.437]$ & 0.773 & 141\,:\,230\,:\,0  & $1.9\!\times\!10^{-31}$ \\
        & 滑块    & 6.677 & $[6.024, 7.342]$ & 1.016 & 229\,:\,140\,:\,2  & $4.3\!\times\!10^{-42}$ \\
      \midrule
      \multirow{3}{*}{GLM 4.6}
        & 3 点    & 0.121 & $[0.081, 0.162]$ & 0.297 & 56\,:\,304\,:\,11  & $4.9\!\times\!10^{-8}$ \\
        & 7 点    & 0.272 & $[0.191, 0.350]$ & 0.353 & 115\,:\,219\,:\,37 & $1.1\!\times\!10^{-10}$ \\
        & 滑块    & 3.148 & $[2.423, 3.892]$ & 0.423 & 126\,:\,223\,:\,22 & $1.5\!\times\!10^{-15}$ \\
      \bottomrule
    \end{tabular}
  \end{table}
\end{frame}

\begin{frame}{研究 1b：27 单元完整统计（2/2）——9 模型全部同方向，27 个 $q$ 均 $<10^{-7}$}
  \scriptsize
  \setlength{\tabcolsep}{3pt}
  \renewcommand{\arraystretch}{0.9}
  \begin{table}
    \centering
    \begin{tabular}{llccccc}
      \toprule
      模型 & 尺度 & $\Delta_{F-M}$ & 95\% CI & $d_z$ & 女\,:\,平\,:\,男 & $q$ (FDR) \\
      \midrule
      \multirow{3}{*}{GPT-5.1}
        & 3 点    & 0.299 & $[0.253, 0.345]$ & 0.653 & 111\,:\,260\,:\,0  & $5.3\!\times\!10^{-25}$ \\
        & 7 点    & 0.458 & $[0.404, 0.512]$ & 0.840 & 167\,:\,201\,:\,3  & $2.7\!\times\!10^{-34}$ \\
        & 滑块    & 5.612 & $[4.814, 6.448]$ & 0.714 & 266\,:\,53\,:\,52  & $8.2\!\times\!10^{-36}$ \\
      \midrule
      \multirow{3}{*}{Kimi K2}
        & 3 点    & 0.288 & $[0.237, 0.337]$ & 0.572 & 116\,:\,246\,:\,9  & $4.8\!\times\!10^{-21}$ \\
        & 7 点    & 0.404 & $[0.321, 0.488]$ & 0.505 & 162\,:\,173\,:\,36 & $1.1\!\times\!10^{-17}$ \\
        & 滑块    & 6.377 & $[5.310, 7.469]$ & 0.593 & 264\,:\,47\,:\,60  & $3.7\!\times\!10^{-31}$ \\
      \midrule
      \multirow{3}{*}{Llama Maverick}
        & 3 点    & 0.140 & $[0.102, 0.181]$ & 0.379 & 55\,:\,313\,:\,3   & $1.3\!\times\!10^{-11}$ \\
        & 7 点    & 0.240 & $[0.189, 0.294]$ & 0.477 & 96\,:\,265\,:\,10  & $1.6\!\times\!10^{-16}$ \\
        & 滑块    & 2.881 & $[2.181, 3.598]$ & 0.409 & 109\,:\,240\,:\,22 & $3.4\!\times\!10^{-13}$ \\
      \midrule
      \multirow{3}{*}{Qwen 2.5 72B}
        & 3 点    & 0.235 & $[0.191, 0.280]$ & 0.553 & 87\,:\,284\,:\,0   & $3.3\!\times\!10^{-20}$ \\
        & 7 点    & 0.256 & $[0.210, 0.302]$ & 0.578 & 96\,:\,274\,:\,1   & $1.2\!\times\!10^{-21}$ \\
        & 滑块    & 6.011 & $[5.108, 6.873]$ & 0.686 & 185\,:\,185\,:\,1  & $1.3\!\times\!10^{-32}$ \\
      \bottomrule
    \end{tabular}
  \end{table}
  \begin{alertblock}{判决}
    \scriptsize \textbf{27/27 单元} $\Delta_{F-M}>0$，$q$ 均 $<10^{-7}$；$d_z \in [0.29, 1.02]$——\textbf{H1b.1 的零假设（语义内容）证伪，H1b.1 本身（评估偏差）支持}。
  \end{alertblock}
\end{frame}

\begin{frame}{研究 1b：9 子图分布一致右偏、三尺度贯通——配对差的右偏是稳健现象}
  \centering
  {\tiny GPT-5.1 \hfill\hfill Qwen 2.5 72B \hfill\hfill DeepSeek V3.2\par}\vspace{-0.1em}
  \begin{minipage}[t]{0.32\textwidth}\centering\includegraphics[width=\linewidth,height=0.20\textheight,keepaspectratio]{\oneb/attack_3pt/panels/fig2_difference_distribution_panel_1.pdf}\end{minipage}\hfill%
  \begin{minipage}[t]{0.32\textwidth}\centering\includegraphics[width=\linewidth,height=0.20\textheight,keepaspectratio]{\oneb/attack_3pt/panels/fig2_difference_distribution_panel_8.pdf}\end{minipage}\hfill%
  \begin{minipage}[t]{0.32\textwidth}\centering\includegraphics[width=\linewidth,height=0.20\textheight,keepaspectratio]{\oneb/attack_3pt/panels/fig2_difference_distribution_panel_6.pdf}\end{minipage}

  \vspace{0.2em}
  \begin{minipage}[t]{0.32\textwidth}\centering\includegraphics[width=\linewidth,height=0.20\textheight,keepaspectratio]{\oneb/attack_7pt_likert/panels/fig2_difference_distribution_panel_1.pdf}\end{minipage}\hfill%
  \begin{minipage}[t]{0.32\textwidth}\centering\includegraphics[width=\linewidth,height=0.20\textheight,keepaspectratio]{\oneb/attack_7pt_likert/panels/fig2_difference_distribution_panel_8.pdf}\end{minipage}\hfill%
  \begin{minipage}[t]{0.32\textwidth}\centering\includegraphics[width=\linewidth,height=0.20\textheight,keepaspectratio]{\oneb/attack_7pt_likert/panels/fig2_difference_distribution_panel_6.pdf}\end{minipage}

  \vspace{0.2em}
  \begin{minipage}[t]{0.32\textwidth}\centering\includegraphics[width=\linewidth,height=0.20\textheight,keepaspectratio]{\oneb/attack_slider_0_100/panels/fig2_difference_distribution_panel_1.pdf}\end{minipage}\hfill%
  \begin{minipage}[t]{0.32\textwidth}\centering\includegraphics[width=\linewidth,height=0.20\textheight,keepaspectratio]{\oneb/attack_slider_0_100/panels/fig2_difference_distribution_panel_8.pdf}\end{minipage}\hfill%
  \begin{minipage}[t]{0.32\textwidth}\centering\includegraphics[width=\linewidth,height=0.20\textheight,keepaspectratio]{\oneb/attack_slider_0_100/panels/fig2_difference_distribution_panel_6.pdf}\end{minipage}

  \vspace{0.2em}
  \footnotesize 行 = 3 点 / 7 点 / 滑块；列 = GPT-5.1 / Qwen 2.5 72B / DeepSeek V3.2。9 子图全部右偏；3 点行仅 3 档、7 点行渐宽、滑块行最连续。
\end{frame}

\begin{frame}{研究 1b：三条跨尺度规律}
  \footnotesize
  \begin{block}{规律 1：哪个模型偏差更大的排名跨尺度相当稳定}
    \small 偏差最大的模型（Gemma 在滑块、GPT-5.1 在 7 点和 3 点）在其他尺度下也排前列；偏差最小的（DeepSeek V3.2、GLM 4.6、Llama Maverick）也类似——下节用 Spearman $\rho_s$ 定量呈现。
  \end{block}
  \begin{block}{规律 2：持平的题项随尺度变细而减少}
    \small 9 模型的中位持平率：$80.1\% \to 62.0\% \to 37.7\%$（3 点 $\to$ 7 点 $\to$ 滑块）。偏差没有消失，只是更细的尺度把更小的差别也记下来了。
  \end{block}
  \begin{block}{规律 3：换成滑块不是对所有模型都会放大 $d_z$}
    \small Gemma（$0.77\to 1.02$）、Claude（$0.69\to 0.75$）、两个 DeepSeek 在滑块下 $d_z$ 更大；但 GPT-5.1（$0.84\to 0.71$）、Kimi K2、Llama 反而持平或略降——"滑块只是分辨率更高的同一套偏差"并不适用于所有模型。
  \end{block}
\end{frame}

\begin{frame}{研究 1b：27/27 单元 $q<.001$、跨尺度强度量级稳健——整体偏差方向三尺度一致}
  \footnotesize
  \setlength{\tabcolsep}{3pt}
  \begin{table}
    \centering
    \begin{tabular}{lcccccc}
      \toprule
      模型 & \multicolumn{2}{c}{3 点} & \multicolumn{2}{c}{7 点 Likert} & \multicolumn{2}{c}{0--100 滑块} \\
      \cmidrule(lr){2-3}\cmidrule(lr){4-5}\cmidrule(lr){6-7}
           & $\Delta_{F-M}$ & $d_z$ & $\Delta_{F-M}$ & $d_z$ & $\Delta_{F-M}$ & $d_z$ \\
      \midrule
      Claude Opus 4.5  & 0.199 & 0.498 & 0.345 & 0.693 & 3.55 & 0.748 \\
      DeepSeek R1      & 0.105 & 0.289 & 0.243 & 0.422 & 4.33 & 0.510 \\
      DeepSeek V3.2    & 0.124 & 0.289 & 0.237 & 0.371 & 3.89 & 0.460 \\
      Gemma 4 31B      & 0.143 & 0.408 & 0.385 & 0.773 & 6.68 & 1.016 \\
      GLM 4.6          & 0.121 & 0.297 & 0.272 & 0.353 & 3.15 & 0.423 \\
      GPT-5.1          & 0.299 & 0.653 & 0.458 & 0.840 & 5.61 & 0.714 \\
      Kimi K2          & 0.288 & 0.572 & 0.404 & 0.505 & 6.38 & 0.593 \\
      Llama 4 Maverick & 0.140 & 0.379 & 0.240 & 0.477 & 2.88 & 0.409 \\
      Qwen 2.5 72B     & 0.235 & 0.553 & 0.256 & 0.578 & 6.01 & 0.686 \\
      \bottomrule
    \end{tabular}
  \end{table}

  \vspace{0.2em}
  \begin{block}{}
    \small 3 尺度下 27/27 单元显著（$q < .001$）；\emph{整体强度}跨尺度稳健、方向全部一致。
  \end{block}
\end{frame}

\begin{frame}{研究 1b：跨模型 $d_z$ 在任意两尺度间贴近对角线——排序稳定，滑块仅稍放大}
  \centering
  \includegraphics[width=0.88\textwidth,height=0.70\textheight,keepaspectratio]{\oneb/placeholder_b_three_scale_dz_scatter.pdf}

  \vspace{0.1em}
  \small
  \begin{itemize}
    \item 每点一模型；三对尺度两两绘制；虚线 $y=x$。
    \item 点贴近对角线 $\Rightarrow$ 跨模型 $d_z$ 排序稳定；滑块下若干模型略位于对角线上方（滑块分辨率高、小偏差被更充分记录）。
  \end{itemize}
\end{frame}

\begin{frame}{研究 1b：9 个模型均把"性化 + 外貌"列打得最深——维度集中性有模型内证据}
  \centering
  \begin{columns}[T,onlytextwidth]
    \column{0.33\textwidth}
      \includegraphics[width=\linewidth,height=0.74\textheight,keepaspectratio]{\oneb/attack_7pt_likert/panels/fig4_level1_forest_panel_1.pdf}\\[-0.2em]
      {\scriptsize GPT-5.1}
    \column{0.33\textwidth}
      \includegraphics[width=\linewidth,height=0.74\textheight,keepaspectratio]{\oneb/attack_7pt_likert/panels/fig4_level1_forest_panel_9.pdf}\\[-0.2em]
      {\scriptsize Gemma 4 31B}
    \column{0.33\textwidth}
      \includegraphics[width=\linewidth,height=0.74\textheight,keepaspectratio]{\oneb/attack_7pt_likert/panels/fig4_level1_forest_panel_6.pdf}\\[-0.2em]
      {\scriptsize DeepSeek V3.2}
  \end{columns}

  \vspace{0.1em}
  \small GPT-5.1 覆盖面最广（7 领域 $d_z \geq 0.77$）；DeepSeek V3.2 最集中（仅 4 领域显著）；性化 + 外貌在所有 9 模型均最强。
\end{frame}

\begin{frame}{研究 1b：列方向（维度）与行方向（模型）共同显示——H1b.2 维度集中性支持}
  \centering
  \includegraphics[width=0.74\textwidth,height=0.66\textheight,keepaspectratio]{\oneb/placeholder_c_model_by_level1_dz_heatmap.pdf}

  \scriptsize \textbf{两条一致性结构}——列方向：性化 + 外貌两列在 9 模型均深红（$d_z$ 多数 $> 0.8$）；行方向：GPT-5.1 与 Gemma 覆盖全行（7 领域 $d_z \geq 0.77$），DeepSeek V3.2 仅少数列显著。
\end{frame}

\begin{frame}{研究 1b：30 单元（9 模型 $\times$ 3 尺度）下"性化 + 外貌"一致深红——H1b.2 稳健}
  \centering
  \includegraphics[width=\textwidth,height=0.80\textheight,keepaspectratio]{\artifacts/model_by_level1_by_scale_dz_heatmap.pdf}

  \vspace{0.1em}
  \footnotesize 9 模型 $\times$ 10 领域 $\times$ 3 尺度的 $d_z$ 热图。\emph{性化 + 外貌两列在三尺度 $\times$ 9 模型的 30 单元里一致深红}；两列合计构成了维度集中性的主要载体——\textbf{H1b.2 支持}。
\end{frame}

\begin{frame}{研究 1b：跨尺度同构的诊断分三层——聚合、维度、题项，自由度递增}
  \footnotesize
  \begin{itemize}
    \item \textbf{聚合层}：每模型 $\to$ 一个 $d_z$ 标量；9 模型构成长 9 向量。
    \item \textbf{维度层}：每模型 $\to$ 10 个一级领域 $d_z$；$n=10$。
    \item \textbf{题项层}：每模型 $\to$ 371 对 $\Delta_i$；$n=371$。
  \end{itemize}

  \vspace{0.5em}
  \begin{block}{}
    \small 分析越细，变量越多，跨尺度的一致性理论上会下降——但在哪一层开始基本不一致，是需要实际数据来回答的。下面依次呈现三层的 Spearman $\rho_s$。
  \end{block}
\end{frame}

\begin{frame}{研究 1b：聚合层 $\rho_s = 0.57$--$0.85$——"谁偏差更强"的排序跨尺度高度一致}
  \centering
  \includegraphics[width=0.84\textwidth,height=0.62\textheight,keepaspectratio]{\artifacts/overall_dz_cross_scale_scatter.pdf}

  \vspace{0.1em}
  \footnotesize
  \setlength{\tabcolsep}{5pt}
  \begin{tabular}{lccc}
    \toprule
    尺度对 & 3pt vs 7pt & 7pt vs slider & 3pt vs slider \\
    \midrule
    Spearman $\rho_s$ & $+0.783$ & $+0.850$ & $+0.567$ \\
    \bottomrule
  \end{tabular}

  \vspace{0.2em}
  \small 聚合层 $\rho_s \geq 0.57$——\emph{哪个模型偏差更大的排名在三种尺度之间相当一致}。
\end{frame}

\begin{frame}{研究 1b：维度层中位 $\rho = 0.35$--$0.48$——模型间异质性明显、部分模型出现负相关}
  \centering
  \includegraphics[width=\textwidth,height=0.60\textheight,keepaspectratio]{\artifacts/dimension_level_cross_scale_rho_forest.pdf}

  \scriptsize 9 模型中位 $\rho$：3pt vs 7pt $=+0.35$，7pt vs slider $=+0.42$，3pt vs slider $=+0.48$——低于聚合层但高于题项层。\textbf{模型间差异明显}：GPT-5.1 与 Gemma 三对均 $>+0.6$；Claude（3pt vs 7pt $\rho=-0.27$）与 DeepSeek V3.2（$\rho=-0.43$）出现负相关——对这些模型而言，"哪几个维度偏差最强"会随尺度变化而改变。
\end{frame}

\begin{frame}{研究 1b：题项层中位 $\rho \approx 0$--$0.15$——"具体哪一句最偏"跨尺度几乎独立}
  \centering
  \includegraphics[width=\textwidth,height=0.56\textheight,keepaspectratio]{\artifacts/within_model_cross_scale_similarity.pdf}

  \footnotesize
  \setlength{\tabcolsep}{5pt}
  \renewcommand{\arraystretch}{0.95}
  \begin{tabular}{lccc}
    \toprule
    层级 & 3pt vs 7pt & 7pt vs slider & 3pt vs slider \\
    \midrule
    题项级中位 $\rho$（9 模型）  & $+0.006$ & $+0.152$ & $+0.116$ \\
    聚合级 $\rho_s$（9 模型 $d_z$） & $+0.783$ & $+0.850$ & $+0.567$ \\
    \bottomrule
  \end{tabular}
\end{frame}

\begin{frame}{研究 1b：H1b.3 的结果因分析层级而异，需分层汇报}
  \footnotesize
  \setlength{\tabcolsep}{4pt}
  \begin{table}
    \centering
    \begin{tabular}{lcccc}
      \toprule
      层级 & 3pt vs 7pt & 7pt vs slider & 3pt vs slider & 判决 \\
      \midrule
      聚合层（$n=9$ 模型 $d_z$） & $+0.78$ & $+0.85$ & $+0.57$ & \textbf{支持} \\
      维度层（$n=10$ 领域 $d_z$，中位） & $+0.35$ & $+0.42$ & $+0.48$ & \textbf{部分支持} \\
      题项层（$n=371$ 对 $\Delta_i$，中位） & $+0.006$ & $+0.152$ & $+0.116$ & \textbf{证伪} \\
      \bottomrule
    \end{tabular}
  \end{table}

  \vspace{0.3em}
  \begin{block}{分层结论}
    \small 简单来说——"三种尺度在测同一套偏差"这个直觉，在整体 $d_z$ 排名上是对的，在维度层面部分成立，但在题项层面不成立。题项级别的个别差异在求均值时被抵消了；维度层面介于两者之间，同时模型间差异最大。
  \end{block}
\end{frame}

\begin{frame}{研究 1b：持平不是中立——偏差被集中到少数决策分化的题项上}
  \centering
  \begin{columns}[T,onlytextwidth]
    \column{0.48\textwidth}
      \includegraphics[width=\linewidth,height=0.34\textheight,keepaspectratio]{\oneb/attack_7pt_likert/panels/fig1_paired_scores_panel_2.pdf}\\[-0.2em]
      {\scriptsize Claude Opus 4.5（持平 65.0\%，不持平时 64:1）}
    \column{0.48\textwidth}
      \includegraphics[width=\linewidth,height=0.34\textheight,keepaspectratio]{\oneb/attack_7pt_likert/panels/fig1_paired_scores_panel_4.pdf}\\[-0.2em]
      {\scriptsize Llama Maverick（持平 71.4\%，不持平时 9.6:1）}
  \end{columns}
  \begin{columns}[T,onlytextwidth]
    \column{0.48\textwidth}
      \includegraphics[width=\linewidth,height=0.34\textheight,keepaspectratio]{\oneb/attack_7pt_likert/panels/fig1_paired_scores_panel_8.pdf}\\[-0.2em]
      {\scriptsize Qwen 2.5 72B（持平 73.9\%，不持平时 96:1）}
    \column{0.48\textwidth}
      \includegraphics[width=\linewidth,height=0.34\textheight,keepaspectratio]{\oneb/attack_7pt_likert/panels/fig1_paired_scores_panel_1.pdf}\\[-0.2em]
      {\scriptsize GPT-5.1（持平 54.2\%，不持平时 55.7:1）}
  \end{columns}

  \vspace{0.1em}
  \small 高持平比例模型一旦打破持平、方向几乎全部偏向女人版得分更高；\emph{持平不是中立}——偏差被集中到少数决策分化的题项上。
\end{frame}

\begin{frame}{研究 1b 小结：H1b.1／H1b.2 支持、H1b.3 分层——整体量级稳、具体题项重组}
  \footnotesize
  \setlength{\tabcolsep}{4pt}
  \begin{table}
    \centering
    \begin{tabular}{lll}
      \toprule
      假设 & 关键预测 & 数据判决 \\
      \midrule
      H1b.1 评估偏差 & $\Delta_{F-M}>0$，9/9 方向一致 & \textbf{支持}（27/27 单元 $q<10^{-7}$；$d_z$ 跨度 $0.29$--$1.02$；零假设证伪） \\
      H1b.2 维度集中 & 性化 + 外貌最强 & \textbf{支持}（两列在 9 模型均显著、$d_z$ 最大） \\
      H1b.3 题项级同构 & 跨尺度 $\rho \approx 1$ & \textbf{题项级证伪}（$\rho$ 中位 $0.006$--$0.15$）；\textbf{聚合级支持}（$\rho_s = 0.57$--$0.85$） \\
      \bottomrule
    \end{tabular}
  \end{table}

  \vspace{0.3em}
  \begin{block}{一行结论}
    \small 9 个 LLM 一致把"女人版"打得更高，偏差集中在身体／性相关维度；整体效应量在三种尺度之间相当稳定，但\emph{具体哪一条语句}偏差最大，换一种尺度结果就不同了。
  \end{block}
\end{frame}

% ===========================================================================
\section{研究 2a：779 个人类评分 session 性别配对评分}

\begin{frame}{研究 2a：研究问题与设计}
  \small\emph{从 1b 到 2a 的过渡}：1b 告诉我们 9 个 LLM 对女人版一致打得更高，但 LLM 的训练语料本身就包含了人类标注——单看 LLM 无法区分"模型在模仿人类偏差"与"模型放大了人类偏差"。2a 把同一套最小配对带回人类被试，取得了可与 LLM 直接比对的基线。

  \begin{block}{核心问题}
    \small 在研究 1b 的同一套 371 对中文配对语料上，招募 712 名中文母语被试打分（合计 779 个有效评分 session，\textbf{全文以 session 为分析单元}）；人类是否与 LLM 同方向？被试性别与尺度谁是更大的异质来源？
  \end{block}

  \begin{itemize}
    \item \textbf{被试规模}：779 个有效评分 session（对应 712 名独立被试：645 人各贡献 1 个 session，67 人各贡献 2 个 session；其中 34 人在两个不同尺度条件下各留下 1 个 session，其余 33 人在同一条件内有两次提交）。
    \item \textbf{条件分配}：$n_{3pt} = 268$、$n_{7pt} = 244$、$n_{slider} = 267$；同一条件内被试间随机分派到女人版或男人版。
    \item \textbf{分析}：每题项、每尺度、每被试子群分别计算 $\Delta_{F-M}$，在 371 对上汇总为 $d_z$。
  \end{itemize}
\end{frame}

\begin{frame}{研究 2a：被试信息（session $\times$ 性别 $\times$ 尺度）}
  \footnotesize
  \setlength{\tabcolsep}{6pt}
  \begin{table}
    \centering
    \caption{\small 779 个有效评分 session 按被试性别与评分尺度的 $2 \times 3$ 交叉分布}
    \begin{tabular}{lccc|c}
      \toprule
      尺度 & 女性被试 & 男性被试 & 合计 & 占比 \\
      \midrule
      3 点        & 134 & 134 & 268 & 34.4\% \\
      7 点 Likert & 128 & 116 & 244 & 31.3\% \\
      0--100 滑块 & 130 & 137 & 267 & 34.3\% \\
      \midrule
      合计        & \textbf{392} & \textbf{387} & \textbf{779} & 100\% \\
      占比        & 50.3\%       & 49.7\%       & 100\%        &        \\
      \bottomrule
    \end{tabular}
  \end{table}

  \vspace{0.3em}
  \begin{itemize}
    \small
    \item \textbf{独立被试人数 $= 712$}；其中 $645$ 人只留下 $1$ 个 session、$67$ 人留下 $2$ 个 session。合计 $645 \times 1 + 67 \times 2 = 779$ session。
    \item 性别在总体上几乎平衡（女 $50.3\%$ / 男 $49.7\%$）；各尺度条件内也近似平衡（女性占比 $49\%$--$53\%$）。
    \item 全文后续 2a 所有 $d_z$、ICC、$\rho$ 均以 \emph{779 个 session} 为分析单元。
  \end{itemize}
\end{frame}

\begin{frame}{研究 2a：被试人口学特征（按尺度条件拆分）}
  \tiny
  \setlength{\tabcolsep}{3pt}
  \renewcommand{\arraystretch}{0.92}
  \begin{table}
    \centering
    \begin{tabular}{lrrrr}
      \toprule
      变量 / 取值 & 3 点（$n{=}268$） & 7 点（$n{=}244$） & 滑块（$n{=}267$） & 全样本（$n{=}779$） \\
      \midrule
      \multicolumn{5}{l}{\textit{年龄段}} \\
      \quad 18--24 岁 & 176 (65.7\%) & 162 (66.4\%) & 200 (74.9\%) & 538 (69.1\%) \\
      \quad 24--34 岁 & 84 (31.3\%)  & 73 (29.9\%)  & 53 (19.9\%)  & 210 (27.0\%) \\
      \quad 35--50 岁 & 6 (2.2\%)    & 6 (2.5\%)    & 12 (4.5\%)   & 24 (3.1\%)   \\
      \quad 其他（12--18／50+） & 2 (0.7\%) & 3 (1.2\%) & 2 (0.7\%) & 7 (0.9\%) \\
      \midrule
      \multicolumn{5}{l}{\textit{学历}} \\
      \quad 本科 & 205 (76.5\%) & 183 (75.0\%) & 210 (78.7\%) & 598 (76.8\%) \\
      \quad 大专 & 30 (11.2\%)  & 35 (14.3\%)  & 28 (10.5\%)  & 93 (11.9\%)  \\
      \quad 硕士 & 23 (8.6\%)   & 18 (7.4\%)   & 19 (7.1\%)   & 60 (7.7\%)   \\
      \quad 博士及以上 & 2 (0.7\%) & 2 (0.8\%) & 1 (0.4\%) & 5 (0.6\%) \\
      \quad 高中及以下 & 6 (2.2\%) & 5 (2.0\%) & 8 (3.0\%) & 19 (2.4\%) \\
      \quad 不便回答 & 2 (0.7\%) & 1 (0.4\%) & 1 (0.4\%) & 4 (0.5\%) \\
      \midrule
      \multicolumn{5}{l}{\textit{月收入}} \\
      \quad 无固定收入 & 48 (17.9\%) & 57 (23.4\%) & 62 (23.2\%) & 167 (21.4\%) \\
      \quad 3000 元及以下 & 47 (17.5\%) & 46 (18.9\%) & 66 (24.7\%) & 159 (20.4\%) \\
      \quad 3001--5000 元 & 39 (14.6\%) & 37 (15.2\%) & 50 (18.7\%) & 126 (16.2\%) \\
      \quad 5001--8000 元 & 81 (30.2\%) & 55 (22.5\%) & 54 (20.2\%) & 190 (24.4\%) \\
      \quad 8001--12000 元 & 32 (11.9\%) & 25 (10.2\%) & 18 (6.7\%) & 75 (9.6\%) \\
      \quad 12000 元以上 & 12 (4.5\%)  & 13 (5.3\%) & 12 (4.5\%) & 37 (4.7\%) \\
      \quad 不便回答 & 9 (3.4\%) & 11 (4.5\%) & 5 (1.9\%) & 25 (3.2\%) \\
      \midrule
      \multicolumn{5}{l}{\textit{是否在校学生}} \\
      \quad 是 & 16 (6.0\%)   & 19 (7.8\%)   & 9 (3.4\%)    & 44 (5.6\%)   \\
      \quad 否 & 252 (94.0\%) & 225 (92.2\%) & 258 (96.6\%) & 735 (94.4\%) \\
      \bottomrule
    \end{tabular}
  \end{table}

  \vspace{0.1em}
  \scriptsize \textbf{结构可比}：3 组在年龄、学历、学生比例上高度对齐；\textbf{样本画像}：18--34 岁青年 $96.1\%$、本科及以上 $85.2\%$、月收入 $\leq 8000$ 元 $82.5\%$——典型中文互联网青年样本；收入分布的波动在随机分派自然范围内，不构成系统性分层。
\end{frame}

\begin{frame}{研究 2a：四条可证伪假设（重排以与 1b 配对）}
  \footnotesize
  \begin{block}{H2a.1 F$>$M 主效应假设（对应 H1b.1）}
    \small 人类被试在三种尺度下均呈 $\Delta_{F-M} > 0$，方向与模型一致——在人类端重复 1b 的发现。
  \end{block}
  \begin{block}{H2a.2 维度集中假设（对应 H1b.2）}
    \small 人类端 10 个一级领域中，$d_z$ 同样集中在"性化攻击"与"外貌形象攻击"，与 H1b.2 的模型结果重合；且被试性别会进一步调节这一结构——女性被试在这两个维度上尤其突出。
  \end{block}
  \begin{block}{H2a.3 被试性别差异假设（2a 特有）}
    \small 女性被试的 $d_z$ 显著大于男性被试（日常经验影响敏感度），两组方向均为 $\Delta_{F-M}>0$。
  \end{block}
  \begin{block}{H2a.4 尺度的影响小于被试性别的影响（2a 特有）}
    \small 在 $2 \times 3$（被试性别 $\times$ 尺度）的设计中，\emph{固定分组、换尺度}时 $d_z$ 变动范围，远小于\emph{固定尺度、换性别分组}时的差距——尺度是测量工具，被试性别才是更重要的变异来源。
  \end{block}
\end{frame}

\begin{frame}{研究 2a：人类三尺度 $d_z = 0.51$--$0.62$——与 LLM 同方向，H2a.1 支持}
  \footnotesize
  \setlength{\tabcolsep}{4pt}
  \begin{table}
    \centering
    \begin{tabular}{lccccc}
      \toprule
      尺度 & $\Delta_{F-M}$ & 95\% CI & $d_z$ & 女\,:\,平\,:\,男 & $q$ (FDR) \\
      \midrule
      3 点        & 0.121 & $[0.101, 0.140]$ & 0.621 & 273 : 3 : 95  & $4.0\times 10^{-26}$ \\
      7 点 Likert & 0.260 & $[0.216, 0.306]$ & 0.588 & 268 : 0 : 103 & $2.8\times 10^{-23}$ \\
      0--100 滑块 & 3.55  & $[2.85, 4.27]$   & 0.512 & 257 : 0 : 114 & $8.2\times 10^{-19}$ \\
      \bottomrule
    \end{tabular}
  \end{table}

  \vspace{0.2em}
  \begin{alertblock}{判决}
    \small 三尺度 $q < 10^{-18}$，$d_z = 0.51$--$0.62$；女人版评分更高的题项占比 $69.3\%$--$73.6\%$——\textbf{H2a.1 支持}，与 LLM 方向一致。
  \end{alertblock}
\end{frame}

\begin{frame}{研究 2a：人类三尺度配对差分布均右偏——与 LLM 端分布形状一致}
  \centering
  \begin{columns}[T,onlytextwidth]
    \column{0.33\textwidth}
      \includegraphics[width=\linewidth,height=0.65\textheight,keepaspectratio]{\humanfig/panels/fig2_difference_distribution_panel_1.pdf}\\[-0.2em]
      {\scriptsize 3 点}
    \column{0.33\textwidth}
      \includegraphics[width=\linewidth,height=0.65\textheight,keepaspectratio]{\humanfig/panels/fig2_difference_distribution_panel_2.pdf}\\[-0.2em]
      {\scriptsize 7 点 Likert}
    \column{0.33\textwidth}
      \includegraphics[width=\linewidth,height=0.65\textheight,keepaspectratio]{\humanfig/panels/fig2_difference_distribution_panel_3.pdf}\\[-0.2em]
      {\scriptsize 0--100 滑块}
  \end{columns}

  \vspace{0.1em}
  \small 三尺度下分布均系统性右偏；虚线为 0；与 9 个 LLM 一致。
\end{frame}

\begin{frame}{研究 2a：女人版更高的题项占比 $69\%$--$74\%$——方向与 9 个 LLM 完全一致}
  \centering
  \includegraphics[width=0.82\textwidth,height=0.72\textheight,keepaspectratio]{\humanfig/fig3_directionality.pdf}

  \vspace{0.1em}
  \small 三尺度下女人版更高的题项占比 $69.3\%$--$73.6\%$；男人版更高仅 $25.6\%$--$30.7\%$——9/9 模型方向与人类方向完全一致。
\end{frame}

\begin{frame}{研究 2a：女性被试 $d_z$ 稳定约为男性被试的 3 倍（3.0／3.6／3.5 倍）}
  \footnotesize
  \setlength{\tabcolsep}{4pt}
  \begin{table}
    \centering
    \begin{tabular}{llccccc}
      \toprule
      尺度 & 被试性别 & $n_{\text{part}}$ & $\Delta_{F-M}$ & 95\% CI & $d_z$ & 女/男 \\
      \midrule
      3 点        & 男 & 134 & 0.059 & $[0.032, 0.086]$ & 0.219 & --- \\
                  & 女 & 134 & 0.180 & $[0.152, 0.207]$ & 0.655 & $3.0\times$ \\
      7 点 Likert & 男 & 116 & 0.126 & $[0.059, 0.190]$ & 0.198 & --- \\
                  & 女 & 128 & 0.393 & $[0.334, 0.449]$ & 0.704 & $3.6\times$ \\
      0--100 滑块 & 男 & 137 & 1.66  & $[0.65, 2.63]$   & 0.167 & --- \\
                  & 女 & 130 & 5.64  & $[4.63, 6.59]$   & 0.590 & $3.5\times$ \\
      \bottomrule
    \end{tabular}
  \end{table}

  \vspace{0.2em}
  \begin{alertblock}{判决}
    \small 女性被试 $d_z$ 稳定约为男性被试的 \textbf{3 倍}（3.0 / 3.6 / 3.5），跨尺度一致——\textbf{H2a.3 支持}。
  \end{alertblock}
\end{frame}

\begin{frame}{研究 2a：女性被试行的右偏显著强于男性——H2a.3 的视觉证据}
  \centering
  \includegraphics[width=\textwidth,height=0.78\textheight,keepaspectratio]{\humanpgfig/fig2_participant_gender_difference_distribution.pdf}

  \vspace{0.1em}
  \small 行 = 被试性别；列 = 评分尺度。\emph{女性被试}行上分布右偏显著强于男性被试行，且该模式跨三尺度一致。
\end{frame}

\begin{frame}{研究 2a：组间差约为组内差的 4--8 倍——被试性别才是主导异质来源}
  \small
  \begin{itemize}
    \item \textbf{固定分组、变动尺度}的组内 $d_z$ 极差：
      \begin{itemize}
        \small
        \item 女性被试：$0.655 \to 0.704 \to 0.590$，极差 $= 0.114$。
        \item 男性被试：$0.219 \to 0.198 \to 0.167$，极差 $= 0.052$。
      \end{itemize}
    \item \textbf{固定尺度、变动分组}的组间 $d_z$ 差：$0.43$--$0.51$。
    \item 后者约为前者的 \textbf{4--8 倍}。
  \end{itemize}

  \vspace{0.3em}
  \begin{alertblock}{判决}
    \small 尺度选择是相对中性的测量决策；被试性别才是决定效应量的关键变量——\textbf{H2a.4 支持}。
  \end{alertblock}
\end{frame}

\begin{frame}{研究 2a：性化 + 外貌恰是 9 个 LLM 最集中的列——人机共同热区}
  \centering
  \includegraphics[width=0.92\textwidth,height=0.78\textheight,keepaspectratio]{\humanpgfig/fig4_participant_gender_level1_forest.pdf}

  \vspace{0.1em}
  \small 人类整体 $d_z$ 集中于\textbf{性化攻击}与\textbf{外貌形象攻击}两列——恰是 1b 所有 9 模型均显著的列（与 H1b.2 重合）；该集中结构在\emph{女性被试}上尤为突出——\textbf{H2a.2 支持}。
\end{frame}

\begin{frame}{研究 2a：人类持平带比 LLM 稀——个体差异大、少有题项平均分完全一致}
  \centering
  \begin{columns}[T,onlytextwidth]
    \column{0.33\textwidth}
      \includegraphics[width=\linewidth,height=0.65\textheight,keepaspectratio]{\humanfig/panels/fig1_paired_scores_panel_1.pdf}\\[-0.2em]
      {\scriptsize 3 点}
    \column{0.33\textwidth}
      \includegraphics[width=\linewidth,height=0.65\textheight,keepaspectratio]{\humanfig/panels/fig1_paired_scores_panel_2.pdf}\\[-0.2em]
      {\scriptsize 7 点 Likert}
    \column{0.33\textwidth}
      \includegraphics[width=\linewidth,height=0.65\textheight,keepaspectratio]{\humanfig/panels/fig1_paired_scores_panel_3.pdf}\\[-0.2em]
      {\scriptsize 0--100 滑块}
  \end{columns}

  \vspace{0.1em}
  \small 横轴 = 男人版题项平均分；纵轴 = 女人版；对角线上方密度 = 女人版更高的题项比例。\emph{持平带显著比 LLM 更稀}——人类评分个体差异大、不易给出完全一致的题项\emph{平均}分。
\end{frame}

\begin{frame}{研究 2a 小结：四条假设全部支持——人机同方向，被试性别是最稳健异质来源}
  \small
  \setlength{\tabcolsep}{4pt}
  \begin{table}
    \centering
    \begin{tabular}{lll}
      \toprule
      假设 & 关键预测 & 数据判决 \\
      \midrule
      H2a.1 主效应延拓（对应 H1b.1） & 人类三尺度 $\Delta_{F-M}>0$ & \textbf{支持}（$d_z = 0.51$--$0.62$，$q<10^{-18}$） \\
      H2a.2 维度集中延拓（对应 H1b.2） & 性化 + 外貌最强，女性被试尤其 & \textbf{支持}（与 LLM 最集中领域重合） \\
      H2a.3 被试性别异质 & 女性 $d_z$ $>$ 男性 $d_z$ & \textbf{支持}（3.0--3.6 倍，跨尺度稳定） \\
      H2a.4 尺度 $\ll$ 性别 & 组内 $\ll$ 组间 & \textbf{支持}（4--8 倍关系） \\
      \bottomrule
    \end{tabular}
  \end{table}

  \vspace{0.3em}
  \begin{block}{一行结论}
    \small 人类与 LLM 在方向、维度集中上一致；被试性别是跨尺度最稳健的异质来源；性化／外貌两域是人机偏差的共同热区。
  \end{block}
\end{frame}

% ===========================================================================
\section{研究 2a$\cdot$ICC（Intraclass Correlation Coefficient，类内相关系数）：人机偏差模式的题项级相似度}

\begin{frame}{研究 2a$\cdot$ICC：研究问题与设计}
  \small\emph{从 2a 到 2a$\cdot$ICC 的过渡}：到这里我们已经确认人机在"整体方向"和"量级排序"上一致；但\textbf{方向一致不等于题项级模式一致}——人机可能在"哪一整类题最严重"上意见相同、却在"具体哪一句最严重"上各走各的。2a$\cdot$ICC（Intraclass Correlation Coefficient，类内相关系数）就是从题项级追问人机对齐。

  \begin{block}{核心问题}
    \small 人机在\emph{方向}与\emph{量级}上已达成一致；在\emph{题项层面}，哪一个 LLM 的偏差分布与人类最接近？\emph{哪一句被判为最偏差}这一问题上人机是否达成共识？
  \end{block}

  \begin{itemize}
    \item \textbf{题项级偏差向量}：$\Delta_i^\text{人} = \bar{y}_{i,\text{女}}^\text{人} - \bar{y}_{i,\text{男}}^\text{人}$；$\Delta_i^\text{模} = y_{i,\text{女}}^\text{模} - y_{i,\text{男}}^\text{模}$（同尺度内比较）。
    \item \textbf{主指标}：ICC(2,1) 绝对一致；对照：ICC(3,1) 一致性。Bootstrap $B = 2000$ 95\% CI。
    \item \textbf{9 个单元}：3 被试子群（全体 / 男 / 女）$\times$ 3 尺度；每单元遍历 9 个 LLM、选 ICC(2,1) 最高者。
  \end{itemize}
\end{frame}

\begin{frame}{研究 2a·ICC：三条假设}
  \footnotesize
  \begin{block}{H2a.ICC.1 单一"最像人"模型假设（零假设）}
    \small 若 9 个 LLM 的对齐已收敛到同一种"类人"结构，则 9 个单元的 top-1 应稳定为同一模型。
  \end{block}
  \begin{block}{H2a.ICC.2 "最像人"模型因子群而异（与 ICC.1 互斥）}
    \small 不同被试子群或尺度下，top-1 会切换到不同的 LLM——每个单元的"最像人"模型各不相同。
  \end{block}
  \begin{block}{H2a.ICC.3 效应量大则题项级也更一致}
    \small $d_z$ 更大的被试子群（女性）应与 LLM 的题项级模式更接近、ICC 更高。若结果相反，说明\emph{整体效应量}与\emph{题项级模式相似度}反映的是两个不同维度的偏差。
  \end{block}
\end{frame}

\begin{frame}{研究 2a·ICC：5 个不同模型在 9 个单元轮流登顶——不存在普适"最像人"模型}
  \footnotesize
  \setlength{\tabcolsep}{4pt}
  \begin{table}
    \centering
    \begin{tabular}{llccccc}
      \toprule
      被试子群 & 尺度 & $n_\text{part}$ & 最相似模型 & ICC(2,1) & 95\% CI & ICC(3,1) \\
      \midrule
      全体 & 3 点        & 268 & Llama 4 Maverick & 0.089 & $[0.003, 0.170]$  & 0.089 \\
      全体 & 7 点 Likert & 244 & Claude Opus 4.5  & 0.158 & $[0.052, 0.264]$  & 0.160 \\
      全体 & 0--100 滑块 & 267 & Claude Opus 4.5  & 0.222 & $[0.067, 0.376]$  & 0.222 \\
      男   & 3 点        & 134 & Llama 4 Maverick & 0.096 & $[-0.001, 0.198]$ & 0.099 \\
      男   & 7 点 Likert & 116 & GLM 4.6          & 0.155 & $[0.056, 0.256]$  & 0.158 \\
      男   & 0--100 滑块 & 137 & DeepSeek R1      & 0.174 & $[0.081, 0.256]$  & 0.180 \\
      女   & 3 点        & 134 & DeepSeek R1      & 0.100 & $[0.021, 0.178]$  & 0.103 \\
      女   & 7 点 Likert & 128 & Claude Opus 4.5  & 0.116 & $[0.010, 0.219]$  & 0.116 \\
      女   & 0--100 滑块 & 130 & Qwen 2.5 72B     & 0.112 & $[-0.008, 0.223]$ & 0.112 \\
      \bottomrule
    \end{tabular}
  \end{table}

  \vspace{0.2em}
  \small \textbf{5 个不同模型}在 9 个单元轮流登顶（Claude $\times 3$、Llama $\times 2$、DeepSeek R1 $\times 2$、GLM $\times 1$、Qwen $\times 1$）。
\end{frame}

\begin{frame}{研究 2a·ICC：Claude 4.5 于 7 点／滑块胜出——ICC 随尺度分辨率上升}
  \centering
  \includegraphics[width=\textwidth,height=0.82\textheight,keepaspectratio]{\artifacts/human_model_delta_similarity.pdf}
\end{frame}

\begin{frame}{研究 2a·ICC：拆开被试性别后 top-1 跳到 GLM／DeepSeek R1／Qwen——不同子群下最相似模型不同}
  \centering
  \includegraphics[width=\textwidth,height=0.82\textheight,keepaspectratio]{\artifacts/human_model_delta_similarity_by_participant_gender.pdf}
\end{frame}

\begin{frame}{研究 2a·ICC：三条主要观察}
  \footnotesize
  \begin{block}{观察一：没有固定的"最像人"模型，子群不同则冠军不同}
    \small 9 个单元里出现 5 个不同 top-1。全体样本下 Claude 4.5 在 7 点与滑块胜出；拆开被试性别后，男性 7 点 = GLM 4.6，女性滑块 = Qwen——因此"某 LLM 最像人"的说法必须注明是哪个被试群体、哪种尺度。
  \end{block}
  \begin{block}{观察二：ICC 随尺度精度提升而上升——3 点尺度下题项级对齐最弱}
    \small 全体被试最大 ICC(2,1) 从 3 点 $0.089$ 升至 7 点 $0.158$、滑块 $0.222$；但 9 个单元全部处于 Koo \& Li (2016) 的 poor 区间（$<0.50$）。3 点尺度只有三档取值，偏差值被压缩，题项间差异难以区分——后续题项级分析建议使用 7 点或滑块。
  \end{block}
  \begin{block}{观察三：效应量大的子群，题项级反而与 LLM 更不像}
    \small 男性被试 7 点／滑块的 CI 均不含 0，即题项级与 LLM 有实质相关；女性被试 CI 反而跨 0。女性被试整体 $d_z$ 更大，却在哪道题最严重这一问题上与 LLM 更不一致——效应量与题项级模式刻画的是不同维度的偏差。据此提出"对齐偏斜"假设。
  \end{block}
\end{frame}

\begin{frame}{研究 2a·ICC 小结：没有固定的"最像人"模型，效应量大不等于题项更一致}
  \small
  \setlength{\tabcolsep}{4pt}
  \begin{table}
    \centering
    \begin{tabular}{lll}
      \toprule
      假设 & 关键预测 & 数据判决 \\
      \midrule
      H2a.ICC.1 单一模型 & 9 单元 top-1 同一模型 & \textbf{证伪}（5 个不同 top-1） \\
      H2a.ICC.2 最像人模型因子群而异 & top-1 随被试子群与尺度切换 & \textbf{支持} \\
      H2a.ICC.3 效应量 $\Rightarrow$ 题项级一致 & $d_z$ 大 $\Rightarrow$ ICC 高 & \textbf{证伪}（$d_z$ 大但 ICC 反而低） \\
      \bottomrule
    \end{tabular}
  \end{table}

  \vspace{0.2em}
  \begin{alertblock}{效应量大不等于题项级一致性高}
    \small 女性被试\emph{整体}效应量大，但\emph{哪道题最严重}与 LLM 的判断重合度较低；男性被试效应量弱，但题项偏好与 LLM 更接近。\\
    \textbf{假设}：LLM 对齐更多捕获"男性评分者视角下的攻击性"（alignment skew）——对齐方向存在偏斜，而非总量不足。可用女性评分者数据做 RLHF 再校准验证。
  \end{alertblock}
\end{frame}

% ===========================================================================
\section{研究 2a$\cdot$CS（Cross-Scale similarity，跨尺度相似度）：人机跨尺度稳健性比较}

\begin{frame}{研究 2a$\cdot$CS：研究问题}
  \small\emph{从 2a$\cdot$ICC 到 2a$\cdot$CS 的过渡}：上一节给出了一个看似令人不安的结果——人机题项级相似度全部落在 poor 档。但这可能有两个解释：一是人机之间本来就不像，二是\textbf{人类自己在三尺度间也不像}。2a$\cdot$CS（Cross-Scale similarity，跨尺度相似度）这一节正是做这个\emph{人类内部自洽}的前置检验，也顺便比较人机谁更稳。

  \begin{block}{核心问题}
    \small 人类被试在三种打分尺度下是否给出相互一致的题项级偏差模式？若不然，与 LLM 相比谁更受尺度影响？
  \end{block}

  \begin{itemize}
    \item 节 2a·ICC 已证明人机题项级 ICC 处于 poor 区间（$<0.5$）；本节进一步追问：\emph{两端各自}在三个尺度间是否\emph{自洽}？
    \item \textbf{题项级}指标：每被试子群在每尺度的 $\Delta_i^\text{人}$（长 371），三对尺度间 Spearman $\rho$，bootstrap CI。
    \item \textbf{聚合级}指标：以 10 个一级领域的 $d_z$ 为单位，三对尺度间 Spearman $\rho_s$——人类侧对照 LLM 侧的 9 模型 $d_z$ 向量（节 1b）。
  \end{itemize}
\end{frame}

\begin{frame}{研究 2a·CS：两条假设}
  \footnotesize
  \begin{block}{H2a.CS.1 人类跨尺度不稳健假设}
    \small 若人类评分易受题面呈现方式、记忆负担与锚点认知影响，则题项级 $\rho$ 应普遍偏低，三对组合中位数大致一致。\\
    \textbf{次级预测}：3pt 与 7pt 同属离散 Likert、应共享更强结构；$\rho_{3pt\text{-}7pt} \gg \rho_{7pt\text{-}slider} \approx \rho_{3pt\text{-}slider}$。
  \end{block}
  \begin{block}{H2a.CS.2 LLM 跨尺度更稳健假设}
    \small 若 LLM 作为"纯语言生成器"在不同尺度下对同一内容给成比例评分（如 3 点给 2、7 点给 $4$--$5$、滑块给 $66$），则 LLM 跨尺度 $\rho$ 应\emph{显著高于}人类。
  \end{block}
\end{frame}

\begin{frame}{研究 2a·CS：人类三个被试群体的题项级跨尺度相关均低于 0.13}
  \centering
  \includegraphics[width=0.88\textwidth,height=0.68\textheight,keepaspectratio]{\artifacts/within_human_cross_scale_similarity.pdf}

  \vspace{0.1em}
  \small 全体 / 男性 / 女性被试在三对尺度组合上均呈 $\rho<0.13$；多数 CI 跨越 $0$——\emph{同一批题，换一种打分尺度，具体哪道题最偏的排序就基本重组}。
\end{frame}

\begin{frame}{研究 2a·CS：人机并列森林——中位 $\rho$ 同量级，模型间分化更大}
  \centering
  \includegraphics[width=\textwidth,height=0.80\textheight,keepaspectratio]{\artifacts/within_rater_cross_scale_similarity_human_vs_llm.pdf}
\end{frame}

\begin{frame}{研究 2a·CS：LLM 与人类跨尺度相关相近，但模型间分化更大——H2a.CS.2 不支持}
  \footnotesize
  \setlength{\tabcolsep}{4pt}
  \begin{table}
    \centering
    \begin{tabular}{lcccc}
      \toprule
      尺度对 & 人类中位 $\rho$ & 人类范围 & LLM 中位 $\rho$ & LLM 范围 \\
      \midrule
      3pt vs 7pt       & 0.105 & $[0.011, 0.118]$ & 0.006 & $[-0.233, +0.293]$ \\
      7pt vs slider    & 0.077 & $[0.008, 0.091]$ & 0.152 & $[-0.017, +0.355]$ \\
      3pt vs slider    & 0.108 & $[0.066, 0.124]$ & 0.116 & $[-0.043, +0.229]$ \\
      \bottomrule
    \end{tabular}
  \end{table}

  \begin{itemize}
    \small
    \item \textbf{人类}：9 个单元 $\rho$ 全部 $< 0.13$——题项级跨尺度显著重组。
    \item \textbf{LLM}：中位与人类同量级；3pt vs 7pt 上 LLM 中位 $0.006$ 甚至 \emph{低于}人类 $0.105$。
    \item \textbf{模型间分化}：LLM 跨模型跨度是人类跨被试子群的 \textbf{4--5 倍}（GPT-5.1／Gemma 稳健，Claude 负 $\rho$ 显著）。
  \end{itemize}
\end{frame}

\begin{frame}{研究 2a·CS：聚合级 $+0.86/+0.72/+0.58$ 复刻 LLM 端排序——人机共识}
  \footnotesize
  \setlength{\tabcolsep}{5pt}
  \begin{table}
    \centering
    \begin{tabular}{llccc}
      \toprule
      评分者 & 分析单位 & 3pt vs 7pt & 7pt vs slider & 3pt vs slider \\
      \midrule
      LLM     & 9 模型整体 $d_z$         & $+0.783$（.013）   & $+0.850$（.004）  & $+0.567$（.112） \\
      \midrule
      全体被试 & 10 领域 $d_z$             & $+0.855$（.002）   & $+0.721$（.019）  & $+0.576$（.082） \\
      男性被试 & 10 领域 $d_z$             & $+0.321$（.37）    & $+0.188$（.60）   & $+0.770$（.009） \\
      女性被试 & 10 领域 $d_z$             & $+0.915$（$<$.001）& $+0.552$（.10）   & $+0.442$（.20）  \\
      \bottomrule
    \end{tabular}
  \end{table}

  \vspace{0.2em}
  \begin{alertblock}{聚合级的共识}
    \small 全体被试 $+0.86 / +0.72 / +0.58$ \textbf{几乎完全复刻} LLM 端的 $+0.78 / +0.85 / +0.57$。
    \\女性被试与 LLM 均指向同一个排序：3pt 与 7pt 之间最相似、3pt 与滑块之间最不像——这与"3pt 和 7pt 同属离散 Likert"的直觉一致。
    \\男性被试的 $d_z$ 整体较小，Spearman 不稳定，不构成反例。
  \end{alertblock}
\end{frame}

\begin{frame}{研究 2a·CS 小结：题项级上人机都受尺度影响；在领域层面人机结果相似}
  \small
  \setlength{\tabcolsep}{3pt}
  \begin{table}
    \centering
    \begin{tabular}{lll}
      \toprule
      假设 & 关键预测 & 数据判决 \\
      \midrule
      H2a.CS.1 主张 & 人类题项级 $\rho$ 普遍偏低 & \textbf{支持}（9 个单元 $\rho<0.13$） \\
      H2a.CS.1 次级 & 3pt-7pt $\gg$ 含 slider 两对 & \textbf{题项级证伪}；\textbf{聚合级支持}（全体 $+0.86$ 最高） \\
      H2a.CS.2     & LLM $\rho \gg$ 人类 $\rho$ & \textbf{证伪}（同量级；3pt-7pt LLM 甚至更低） \\
      \bottomrule
    \end{tabular}
  \end{table}

  \vspace{0.3em}
  \begin{block}{分层结论}
    \small \textbf{题项级}：人与 LLM 同样受尺度影响，不存在"LLM 更理性客观"。人机之间真正的差别在于\emph{模型间的分散程度}——LLM 跨模型的跨度是人类跨被试子群的 4--5 倍，模型间两极分化明显（GPT-5.1 正相关、Claude 负相关）。\\
    \textbf{聚合级（领域维度）}：人机共同支持"3pt--7pt 最像、3pt--slider 最不像"的排序。
  \end{block}
\end{frame}

% ===========================================================================
\section{综合结论}

\begin{frame}{全文假设清单（1/2）}
  \scriptsize
  \setlength{\tabcolsep}{4pt}
  \begin{table}
    \centering
    \begin{tabular}{lll}
      \toprule
      假设 & 核心预测 & 判决 \\
      \midrule
      \multicolumn{3}{l}{\textit{研究 1a}} \\
      H1a.1 任务定义分歧 & 等方差、排序稳定 & \textbf{证伪} \\
      H1a.2 语言覆盖不均 & 中文方差 $\gg$ 英文方差 & \textbf{支持} \\
      \midrule
      \multicolumn{3}{l}{\textit{研究 1b}（与 2a 配对）} \\
      H1b.1 评估偏差 & $\Delta_{F-M} > 0$，9/9 方向一致 & \textbf{支持}（零假设亦被证伪） \\
      H1b.2 维度集中 & 性化 + 外貌最强 & \textbf{支持} \\
      H1b.3 题项级跨尺度同构 & 跨尺度 $\rho \approx 1$ & \textbf{题项级证伪；聚合级支持} \\
      \midrule
      \multicolumn{3}{l}{\textit{研究 2a}（与 1b 配对）} \\
      H2a.1 主效应延拓 $\leftrightarrow$ H1b.1 & 人类三尺度 $\Delta_{F-M}>0$ & \textbf{支持} \\
      H2a.2 维度集中延拓 $\leftrightarrow$ H1b.2 & 性化 + 外貌最强，女性尤其 & \textbf{支持} \\
      H2a.3 被试性别异质 & 女性 $d_z$ $>$ 男性 $d_z$ & \textbf{支持}（3 倍） \\
      H2a.4 尺度 $\ll$ 性别 & 组内 $\ll$ 组间 & \textbf{支持}（4--8 倍） \\
      \bottomrule
    \end{tabular}
  \end{table}
\end{frame}

\begin{frame}{全文假设清单（2/2）}
  \scriptsize
  \setlength{\tabcolsep}{4pt}
  \begin{table}
    \centering
    \begin{tabular}{lll}
      \toprule
      假设 & 核心预测 & 判决 \\
      \midrule
      \multicolumn{3}{l}{\textit{研究 2a·ICC}} \\
      H2a.ICC.1 单一"最像人" & 9 单元 top-1 同一模型 & \textbf{证伪}（5 个不同 top-1） \\
      H2a.ICC.2 最像人模型因子群而异 & top-1 随被试子群与尺度切换 & \textbf{支持} \\
      H2a.ICC.3 效应量大则题项级也更一致 & $d_z$ 大 $\Rightarrow$ ICC 高 & \textbf{证伪}（女性 $d_z$ 大、ICC 反而更低，提示对齐偏斜） \\
      \midrule
      \multicolumn{3}{l}{\textit{研究 2a·CS}（H2a.CS.1 与 1b 的 H1b.3 题项层配对）} \\
      H2a.CS.1 主张（$\leftrightarrow$ H1b.3 题项层） & 人类题项级 $\rho$ 偏低 & \textbf{支持} \\
      H2a.CS.1 次级（3pt--7pt 最相似） & 3pt-7pt $\gg$ 含 slider 两对 & \textbf{题项级证伪；聚合级支持} \\
      H2a.CS.2 LLM 更稳健 & LLM $\rho \gg$ 人类 $\rho$ & \textbf{证伪}（同量级） \\
      \bottomrule
    \end{tabular}
  \end{table}
\end{frame}

\begin{frame}{全文一行总结}
  \begin{itemize}
    \item \textbf{1a}：中英性能差距主要不来自标注定义分歧，而来自训练语料对中文的不均衡暴露（方差比 $12.7$--$15.0$ 倍）。
    \item \textbf{1b}：9 个 LLM 一致把"女人版"打得更高，集中在身体／性相关维度；\emph{整体量级}跨尺度稳定，但\emph{具体哪一句}跨尺度显著重组。
    \item \textbf{2a}：人类也同方向偏差；女性被试约为男性被试的 \textbf{3 倍}，跨尺度保持；3 倍差距在性化／外貌两域最强。
    \item \textbf{2a·ICC}：没有哪个 LLM 在所有被试子群和尺度下都最像人；女性被试效应量更大，但在哪道题最严重上与 LLM 反而更不一致——模型的对齐方向更偏向男性评分者，而非对齐的总量不够。
    \item \textbf{2a·CS}：题项级上人与 LLM 同样受尺度影响；领域级聚合上两者均支持"3pt--7pt 最像、3pt--slider 最不像"。
  \end{itemize}
\end{frame}

\begin{frame}{方法论推论}
  \begin{itemize}
    \item \textbf{题项级对齐研究应锚定 7 点或滑块}：3 点尺度将 $\Delta_i$ 压缩到至多三档取值，可检测的题项级信号显著衰减；聚合级另当别论。
    \item \textbf{人机相似度声明必须显式交代"人是谁"}：全体样本下 Claude 4.5 领先，但拆开男／女被试后冠军跳到 GLM 4.6、DeepSeek R1、Qwen 2.5 72B——"哪个 LLM 最像人"无跨尺度、跨被试性别的稳定答案。
    \item \textbf{"量级"与"模式"是两个独立维度}：女性被试 $d_z$ 更大但 ICC 更弱——两者可相互偏离。
    \item \textbf{聚合级与题项级须分层陈述}：偏差在聚合上共享、在题项级分化；任何单一层级声明都可能误读。
  \end{itemize}
\end{frame}

\begin{frame}{研究局限与后续方向}
  \begin{itemize}
    \item \textbf{语言范围}：配对分析仅在中文语境下展开；英文平行语料尚未完成，H1b.1 在英文下的表现不做声明。
    \item \textbf{ICC 绝对值}：全部 9 单元 ICC 位于 Koo \& Li (2016) 的 poor 区间；这是"人机题项级偏差排序仅存在弱对齐"的真实度量，而非实验失败。
    \item \textbf{对齐偏斜的直接检验}：以女性评分者数据做 RLHF 再校准，看 ICC 是否系统性上升——可作为后续实验。
    \item \textbf{模型生态}：9 个 LLM 覆盖主流中、西方厂商；拓展到 30B 以下或专用审核模型仍需验证。
    \item \textbf{维度细粒度}：一级领域已覆盖；二级 / 三级维度的细粒度交互尚未开展。
  \end{itemize}
\end{frame}

\begin{frame}{Conclusions：三个带回家的信息}
  \small
  \begin{enumerate}
    \item \textbf{中英性能差距的"方差侧"比"均值侧"更能说明问题}——方差比 $12.7$--$15.0$ 倍、Bartlett / Levene 一致拒绝等方差；这是 H1a.2（语料暴露不均）而非 H1a.1（定义分歧）的诊断性特征。
    \item \textbf{9 个 LLM 与人类被试同方向}地把"女人版"打得更高，共同集中于\emph{性化 + 外貌}两域；整体量级跨尺度稳定、但\emph{具体哪一句最偏}跨尺度显著重组——偏差是一个必须分层陈述的现象。
    \item \textbf{换尺度对效应量的影响，远小于换被试性别的影响}（前者极差约 $0.05$--$0.11$，后者差距 $0.43$--$0.51$，后者约为前者的 4--8 倍）。女性被试 $d_z$ 更大，但题项级与 LLM 的一致性反而更低——模型更能拟合男性评分者识别出的攻击性题项，而不是女性评分者；这种方向性的偏差可称为\textbf{对齐偏斜}（alignment skew），有别于整体对齐不足。
  \end{enumerate}

  \vspace{0.3em}
  \begin{alertblock}{一句话}
    \small 人机的共同偏差落在同一片"身体／性"领域，但究竟\emph{偏哪一句}跨评分者／跨尺度不稳；"LLM 最像人"的回答必须显式写明——\emph{像哪一群人、在哪一个尺度}。
  \end{alertblock}
\end{frame}

\begin{frame}{致谢}
  \centering
  \vspace{1em}
  \Large 感谢聆听，欢迎提问与讨论。\\[0.5em]
  \normalsize 鲍璇\\ \today
\end{frame}

\end{document}
